% Define a global usable date. Must come before StyleTut
\newcommand{\mydate}{02.12.2016}

% Comment/uncomment this line to toggle handout mode
\newcommand{\handout}{}

% Toggle Handout mode by including this line before including style_tut:
% \newcommand{\handout}{}

%% Beamer-Klasse im korrekten Modus
\ifdefined \handout
\documentclass[handout]{beamer} % Handout mode
\else
\documentclass{beamer}
\fi


% Das ist der KIT-Stil
%\usepackage{../TutTexbib/beamerthemekit}
\usepackage[deutsch,titlepage0]{../TutTexbib/KIT/beamerthemeKITmod}
\TitleImage[width=\titleimagewd]{../figures/titlepage.jpg}
%\usetheme[deutsch,titlepage0]{KIT}

% Include PDFs
\usepackage{pdfpages}

% Libertine font (Original GBI font)
\usepackage{libertine}
%\renewcommand*\familydefault{\sfdefault}  %% Only if the base font of the document is to be sans serif

% Nicer math symbols
\usepackage{eulervm}
%\usepackage{mathpazo}

%% Deutsche Silbentrennung und Beschriftungen
\usepackage[ngerman]{babel}

%% UTF-8-Encoding
\usepackage[utf8]{inputenc}

\usepackage{csquotes}

%% Bibliotheken für viele mathematische Symbole
\usepackage{amsmath, amsfonts, amssymb}

%% Anzeigetiefe für Inhaltsverzeichnis: 1 Stufe
\setcounter{tocdepth}{1}

%% Schönere Schriften
\usepackage[TS1,T1]{fontenc}

%% Tabellen
\usepackage{array}
\usepackage{multicol}

%% Bibliothek für Graphiken
\usepackage{graphicx}

%% der wird sowieso in jeder Datei gesetzt
%%\graphicspath{{../figures/}}

%% Hyperlinks
\usepackage{hyperref}
% I don't know why, but this works and only includes sections and NOT subsections in the pdf-bookmarks.
\hypersetup{bookmarksdepth=subsection} 

%\usepackage{lmodern}
\usepackage{colortbl}
\usepackage[absolute,overlay]{textpos}
\usepackage{listings}
\usepackage{forloop}
%\usepackage{algorithmic} % PseudoCode package 

\usepackage{tikz}
\usetikzlibrary{matrix}
\usetikzlibrary{arrows.meta}
\usetikzlibrary{automata}
\usetikzlibrary{tikzmark}

% Needed for gbi-macros
\usepackage{xspace}

%%%%%%%%%%%% INHALT %%%%%%%%%%%%%%%%

%% Wochennummer
\newcounter{weeknum}

%% Titelinformationen
\title[GBI Tutorium, Woche \theweeknum]{Grundbegriffe der Informatik \\ Tutorium 36}
\subtitle{Termin \theweeknum \ | \mydate \\ Thassilo Helmold}
\author[Thassilo Helmold]{Thassilo Helmold}
\institute{KIT -- Karlsruher Institut für Technologie}
\date{\mydate}

%% Titel einfügen
\newcommand{\titleframe}{\frame{\titlepage}}

%% Alles starten mit \starttut{X}
\newcommand{\starttut}[1]{\setcounter{weeknum}{#1}\titleframe\frame{\frametitle{Inhalt}\tableofcontents} \AtBeginSection[]{%
\begin{frame}
	\tableofcontents[currentsection]
\end{frame}\addtocounter{framenumber}{-1}}}


\newcommand{\framePrevEpisode}{
	\begin{frame}
		\centering
		\textbf{In the previous episode of GBI...}
	\end{frame}
}

%% Legacy: Lastframe. Not for further usage!
\newcommand{\lastframe}[4]{\xkcdframe{#1}{#2}{#3}{#4}{}}
\newcommand{\xkcdframe}[5]{
	\frame[plain]{
		\vspace{-#2pt}
		\begin{figure}[H]
			\centering
			\includegraphics[scale=#1]{#3}
			\vspace{-7pt}
			\caption{ \texttt{\url{#4}} }
			\vspace{5pt}
			#5
		\end{figure} 
	}
}

%% Wörter
\newcommand{\code}[1]{$\mathbf{#1}$}

%% Sterne

\newcounter{starsc}
\newcommand{\stars}[1]{
	\hfill
	\begin{minipage}{100px}
		\forloop{starsc}{0}{\value{starsc} < #1}%
		{%
			\includegraphics[scale=0.05]{star-full.pdf} \hspace*{1px}
		}%
		\forloop{starsc}{\value{starsc}}{\value{starsc} < 5}%
		{%
			\includegraphics[scale=0.05]{star-empty.pdf} \hspace*{1px}
		}
		\vspace*{2px}
	\end{minipage}
}

% No stars for me...
\renewcommand{\stars}[1]{}

\newcommand{\slideThanks}{
	\begin{frame}
		\frametitle{Danksagung}
		\begin{block}{}
			Dieser Foliensatz basiert in Teilen auf Folien von:\\[1em]
			Philipp Basler \\
			Nils Braun \\
			Dominik Doerner \\
			Ou Yue \\
		\end{block}
	\end{frame}
}

%% Verbatim
\usepackage{moreverb}
% % \bigtimes abgeschrieben von http://tex.stackexchange.com/questions/14386/importing-a-single-symbol-from-a-different-font
% \DeclareFontFamily{U}{mathx}{\hyphenchar\font45}
% \DeclareFontShape{U}{mathx}{m}{n}{
%       <5> <6> <7> <8> <9> <10> gen * mathx
%       <10.95> mathx10 <12> <14.4> <17.28> <20.74> <24.88> mathx12
%       }{}
% \DeclareSymbolFont{mathx}{U}{mathx}{m}{n}
% \DeclareMathSymbol{\bigtimes}{\mathop}{mathx}{161}

\RequirePackage{xcolor}

\def\9{\square}
%\def\9{\blank}

% f"ur Aussagenlogik
\colorlet{alcolor}{blue}
\RequirePackage{tikz}
\usetikzlibrary{arrows.meta}
\newcommand{\alimpl}{\mathrel{\tikz[x={(0.1ex,0ex)},y={(0ex,0.1ex)},>={Classical TikZ Rightarrow[]}]{\draw[alcolor,->,line width=0.7pt,line cap=round] (0,0) -- (15,0);\path (0,-6);}}}
\newcommand{\aleqv}{\mathrel{\tikz[x={(0.1ex,0ex)},y={(0ex,0.1ex)},>={Classical TikZ Rightarrow[]}]{\draw[alcolor,<->,line width=0.7pt,line cap=round] (0,0) -- (18,0);\path (0,-6);}}}
\newcommand{\aland}{\mathbin{\raisebox{-0.6pt}{\rotatebox{90}{\texttt{\color{alcolor}\char62}}}}}
\newcommand{\alor}{\mathbin{\raisebox{-0.8pt}{\rotatebox{90}{\texttt{\color{alcolor}\char60}}}}}
%\newcommand{\ali}[1]{_{\mathtt{\color{alcolor}#1}}}
\newcommand{\alv}[1]{\mathtt{\color{alcolor}#1}}
\newcommand{\alnot}{\mathop{\tikz[x={(0.1ex,0ex)},y={(0ex,0.1ex)}]{\draw[alcolor,line width=0.7pt,line cap=round,line join=round] (0,0) -- (10,0) -- (10,-4);\path (0,-8) ;}}}
\newcommand{\alP}{\alv{P}} %ali{#1}}
%\newcommand{\alka}{\negthinspace\hbox{\texttt{\color{alcolor}(}}}
\newcommand{\alka}{\negthinspace\text{\texttt{\color{alcolor}(}}}
%\newcommand{\alkz}{\texttt{\color{alcolor})}}\negthinspace}
\newcommand{\alkz}{\text{\texttt{\color{alcolor})}}\negthinspace}
\newcommand{\AAL}{A_{AL}}
\newcommand{\LAL}{\hbox{\textit{For}}_{AL}}
\newcommand{\AxAL}{\hbox{\textit{Ax}}_{AL}}
\newcommand{\AxEq}{\hbox{\textit{Ax}}_{Eq}}
\newcommand{\AxPL}{\hbox{\textit{Ax}}_{PL}}
\newcommand{\AALV}{\hbox{\textit{Var}}_{AL}}
\newcommand{\MP}{\hbox{\textit{MP}}}
\newcommand{\GEN}{\hbox{\textit{GEN}}}
\newcommand{\W}{\ensuremath{\hbox{\textbf{w}}}\xspace}
\newcommand{\F}{\ensuremath{\hbox{\textbf{f}}}\xspace}
\newcommand{\WF}{\ensuremath{\{\W,\F\}}\xspace}
\newcommand{\val}{\hbox{\textit{val}}}
\newcommand{\valDIb}{\val_{D,I,\beta}}

\newcommand*{\from}{\colon}

% die nachfolgenden Sachen angepasst an cmtt
\newlength{\ttquantwd}
\setlength{\ttquantwd}{1ex}
\newlength{\ttquantht}
\setlength{\ttquantht}{6.75pt}
\def\plall{%
  \tikz[line width=0.67pt,line cap=round,line join=round,baseline=(B),alcolor] {
    \draw (-0.5\ttquantwd,\ttquantht) -- node[coordinate,pos=0.4] (lll){} (-0.25pt,-0.0pt) -- (0.25pt,-0.0pt) -- node[coordinate,pos=0.6] (rrr){} (0.5\ttquantwd,\ttquantht);
    \draw (lll) -- (rrr);
    \coordinate (B) at (0,-0.35pt);
  }%
}
\def\plexist{%
  \tikz[line width=0.67pt,line cap=round,line join=round,baseline=(B),alcolor] {
    \draw (-0.9\ttquantwd,\ttquantht) -- (0,\ttquantht) -- node[coordinate,pos=0.5] (mmm){} (0,0) --  (-0.9\ttquantwd,0);
    \draw (mmm) -- ++(-0.75\ttquantwd,0);
    \coordinate (B) at (0,-0.35pt);
  }\ensuremath{\,}%
}
\let\plexists=\plexist
\newcommand{\NT}[1]{\ensuremath{\langle\mathrm{#1} \rangle}}

\newcommand{\CPL}{\text{\itshape Const}_{PL}}
\newcommand{\FPL}{\text{\itshape Fun}_{PL}}
\newcommand{\RPL}{\text{\itshape Rel}_{PL}}
\newcommand{\VPL}{\text{\itshape Var}_{PL}}
\newcommand{\ATer}{A_{\text{\itshape Ter}}}
\newcommand{\ARel}{A_{\text{\itshape Rel}}}
\newcommand{\AFor}{A_{\text{\itshape For}}}
\newcommand{\LTer}{L_{\text{\itshape Ter}}}
\newcommand{\LRel}{L_{\text{\itshape Rel}}}
\newcommand{\LFor}{L_{\text{\itshape For}}}
\newcommand{\NTer}{N_{\text{\itshape Ter}}}
\newcommand{\NRel}{N_{\text{\itshape Rel}}}
\newcommand{\NFor}{N_{\text{\itshape For}}}
\newcommand{\PTer}{P_{\text{\itshape Ter}}}
\newcommand{\PRel}{P_{\text{\itshape Rel}}}
\newcommand{\PFor}{P_{\text{\itshape For}}}

\newcommand{\plka}{\alka}
\newcommand{\plkz}{\alkz}
%\newcommand{\plka}{\plfoo{(}}
%\newcommand{\plkz}{\plfoo{)}}
\newcommand{\plcomma}{\hbox{\texttt{\color{alcolor},}}}
\newcommand{\pleq}{{\color{alcolor}\,\dot=\,}}

% MODIFIED (DJ)
% previously: \newcommand{\plfoo}[1]{\mathtt{\color{alcolor}#1}}
\newcommand{\plfoo}[1]{\texttt{\color{alcolor}#1}}

\newcommand{\plc}{\plfoo{c}}
\newcommand{\pld}{\plfoo{d}}
\newcommand{\plf}{\plfoo{f}}
\newcommand{\plg}{\plfoo{g}}
\newcommand{\plh}{\plfoo{h}}
\newcommand{\plx}{\plfoo{x}}
\newcommand{\ply}{\plfoo{y}}
\newcommand{\plz}{\plfoo{z}}
\newcommand{\plR}{\plfoo{R}}
\newcommand{\plS}{\plfoo{S}}

\newcommand{\bv}{\mathrm{bv}}
\newcommand{\fv}{\mathrm{fv}}

%\newcommand{\AxAL}{\hbox{\textit{Ax}}_{AL}}
%\newcommand{\AALV}{\hbox{\textit{Var}}_{AL}}

%\renewcommand{\#}[1]{\literal{#1}}
\newcommand{\A}{\mathcal{A}}
\newcommand{\Adr}{\text{Adr}}
\newcommand{\ar}{\mathrm{ar}}
\newcommand{\ascii}[1]{\literal{\char#1}}
%\newcommand{\assert}[1]{\text{/\!\!/\ } #1}
\newcommand{\assert}[1]{\colorbox{black!7!white}{\ensuremath{\{\;#1\;\}}}}
\newcommand{\Assert}[1]{$\langle$\textit{#1}$\rangle$}
\newcommand{\B}{\mathcal{B}}
\newcommand{\bfmod}{\mathbin{\kw{ mod }}}
\newcommand{\bb}{{\text{bb}}}
\def\bottom{\hbox{\small$\pmb{\bot}$}}
\newcommand{\card}[1]{|#1|}
%\newcommand{\cod}{\mathop{\text{cod}}}  % ist in thwmathabbrevs
\newcommand{\Conf}{\mathcal{C}}
\newcommand{\define}[1]{\emph{#1}}
%\renewcommand{\dh}{d.\,h.\@\xspace}
%\newcommand{\Dh}{D.\,h.\@\xspace}
%\newcommand{\engl}[1]{engl.\xspace\emph{#1}}
\newcommand{\eps}{\varepsilon}
%\newcommand{\evtl}{evtl.\@\xspace}
\newcommand{\fbin}{\text{bin}}
\newcommand{\finv}{\text{inv}}
\newcommand{\fnum}{\text{num}}
\newcommand{\fNum}{{\text{Num}}}
\newcommand{\frepr}{\text{repr}}
\newcommand{\fRepr}{\text{Repr}}
\newcommand{\fZkpl}{\text{Zkpl}}
\newcommand{\fLen}{\text{Len}}
\newcommand{\fsem}{\text{sem}}
\providecommand{\fspace}{\mathord{\text{space}}}
\providecommand{\fSpace}{\mathord{\text{Space}}}
\providecommand{\ftime}{\mathord{\text{time}}}
\providecommand{\fTime}{\mathord{\text{Time}}}
\newcommand{\fTrans}{\text{Trans}}
\newcommand{\fVal}{\text{Val}}

% MODIFIED (DJ)
\newcommand{\Val}{\text{Val}}

%\def\G{\mathbb{Z}}
\newcommand{\HT}[1]{\normalfont\textsc{HT-#1}}
\newcommand{\htr}[3]{\{#1\}\;#2\; \{#3\}}
\newcommand{\Id}{\text{I}}
%\newcommand{\ie}{i.\,e.\@\xspace}
\newcommand{\instr}[2]{\texttt{#1}\ \textit{#2}}
\newcommand{\Instr}[2]{\texttt{#1}\ \textrm{#2}}
\newcommand{\instrr}[3]{\texttt{#1}\ \textit{#2}\texttt{(#3)}}
\newcommand{\Instrr}[3]{\texttt{#1}\ \textrm{#2}\texttt{(#3)}}
\newcommand{\io}{\!\mid\!}
\usepackage{KITcolors}
\newcommand{\literal}[1]{\hbox{\textcolor{blue!95!white}{\textup{\texttt{\scalebox{1.11}{#1}}}}}}
%\newcommand{\literal}[1]{\hbox{\textcolor{KITblue!80!black}{\textup{\texttt{#1}}}}}
\def\kasten#1{\leavevmode\literal{\setlength{\fboxsep}{1pt}\fbox{\vrule  width 0pt height 1.5ex depth 0.5ex #1}}}
\newcommand{\kw}[1]{\ensuremath{\mathbf{#1}}}
\newcommand{\lang}[1]{\ensuremath{\langle#1\rangle}}
%\newcommand{\maw}{m.\,a.\,w.\@\xspace}
%\newcommand{\MaW}{M.\,a.\,w.\@\xspace}
\newcommand{\mdefine}[2][FOOBAR]{\define{#2}\def\foobar{FOOBAR}\def\optarg{#1}\ifx\foobar\optarg\def\optarg{#2}\fi\graffito{\optarg}}
\newcommand{\meins}{\rotatebox[origin=c]{180}{1}}
\newcommand{\Mem}{\text{Mem}}
\newcommand{\memread}{\text{memread}}
\newcommand{\memwrite}{\text{memwrite}}
\providecommand{\meta}[1]{\ensuremath{\langle}\textit{#1}\ensuremath{\rangle}}
%\newcommand{\N}{\mathbb{N}}
\newcommand{\NP}{\mathbf{NP}}
\newcommand{\Nadd}{N_{\text{add}}}
\newcommand{\Nmult}{N_{\text{mult}}}
% MODIFIED (DJ): added \!, mathcal{O}
\newcommand{\Oh}[1]{\mathcal{O}\!\left(#1\right)}
\newcommand{\Om}[1]{\Omega\!\left(#1\right)}
\newcommand{\personname}[1]{\textsc{#1}}
\newcommand{\regname}[1]{\texttt{#1}}
\newcommand{\mima}{\textsc{Mima}\xspace}
\newcommand{\mimax}{\textsc{Mima-X}\xspace}

\def\Pclass{\text{\bfseries P}}
\def\PSPACE{\text{\bfseries PSPACE}}

\newcommand{\SPush}{\text{push}}
\newcommand{\SPop}{\text{pop}}
\newcommand{\SPeek}{\text{peek}}
\newcommand{\STop}{\text{top}}
\newcommand{\STos}{\text{\itshape tos}}
\newcommand{\SBos}{\text{\itshape bos}}

%\newcommand{\R}{\mathbb{R}}
\newcommand{\Rnullplus}{\R_0^{+}}
\newcommand{\Rplus}{\R_{+}}
\newcommand{\resp}{resp.\@\xspace}
\newcommand{\Sem}{\text{Sem}}
\newcommand{\sgn}{\mathop{\text{sgn}}}
\newcommand{\sqbox}{\mathop{\raisebox{-6.2pt}{\hbox{\hbox to 0pt{$^{^{\sqcap}}$\hss}$^{^{\sqcup}}$}}}}
\newcommand{\sqleq}{\sqsubseteq}
\newcommand{\sqgeq}{\sqsupseteq}
% MODIFIED (DJ): added \!
\newcommand{\Th}[1]{\Theta\!\left(#1\right)}
%\newcommand{\usw}{usw.\@\xspace}
\newcommand{\V}[1]{\hbox{\textit{#1}}}
\newcommand{\x}{\times}
\newcommand{\ZK}{\mathbb{K}}
%\newcommand{\Z}{\mathbb{Z}}
\newcommand{\zB}{z.\,B.\@\xspace}
\newcommand{\ZB}{Z.\,B.\@\xspace}
% \newcommand{\bb}{{\text{bb}}}
% \def\##1{\hbox{\textcolor{darkblue}{\texttt{#1}}}}
% \def\A{\mathcal{A}}
% \newcommand{\0}{\#0}
% \newcommand{\1}{\#1}
% \newcommand{\Obj}{\text{Obj}}
% \newcommand{\start}{\mathop{\text{start}}}
% \newcommand{\compactlist}{\addtolength{\itemsep}{-\parskip}}
% \newcommand{\fval}{\text{val}}
% \newcommand{\lang}[1]{\ensuremath{\langle#1\rangle}}
% \newcommand{\io}{\!\mid\!}
% \def\sqbox{\mathop{\raisebox{-6.2pt}{\hbox{\hbox to 0pt{$^{^{\sqcap}}$\hss}$^{^{\sqcup}}$}}}}
% \def\sqleq{\sqsubseteq}
% \def\sqgeq{\sqsupseteq}
\def\Td{T_{\overline{d}}}
% \newcommand{\csym}[1]{\ensuremath{\#{c}_{\#{\hbox{\scriptsize #1}}}}}
% \newcommand{\F}{\ensuremath{\mathcal{F}}}
% \newcommand{\fsym}[2]{\ensuremath{\#{f}^{\#{\hbox{\scriptsize #1}}}_{\#{\hbox{\scriptsize #2}}}}}
% \newcommand{\rsym}[2]{\ensuremath{\#{R}^{\#{\hbox{\scriptsize #1}}}_{\#{\hbox{\scriptsize #2}}}}}
% \newcommand{\xsym}[1]{\ensuremath{\#{x}_{\#{\hbox{\scriptsize #1}}}}}
% \newcommand{\I}{\mathcal{I}}
% ********************************************************************

\usepackage{../TutTexbib/thwregex}

\newcommand{\nM}{\mathbb{M}}
\newcommand{\nR}{\mathbb{R}}
\newcommand{\nN}{\mathbb{N}}
\newcommand{\nZ}{\mathbb{Z}}
\newcommand{\nQ}{\mathbb{Q}}
\newcommand{\nB}{\mathbb{B}}
\newcommand{\nC}{\mathbb{C}}
\newcommand{\nK}{\mathbb{K}}
\newcommand{\nF}{\mathbb{F}}
\newcommand{\nG}{\mathbb{G}}
\newcommand{\nullel}{\mathcal{O}}
\newcommand{\einsel}{\mathds{1}}
\newcommand{\nP}{\mathbb{P}}
\newcommand{\Pot}{\mathcal{P}}
\renewcommand{\O}{\text{O}}

%\newcommand{\literal}[1]{\textcolor{blue}{\texttt{#1}}}
\newcommand{\realTilde}{\textasciitilde \ }
\newcommand{\setsize}[1]{\; \mid #1 \mid \; }

\definecolor{myRed}{RGB}{255,75,20}
\colorlet{myGreen}{KITpalegreen}

\newcounter{tfqtempcount}
\newcommand{\truefalseQ}[4]{
	\setcounter{tfqtempcount}{#1}
	\addtocounter{tfqtempcount}{1}
	\truefalseQuestion{#1}{\value{tfqtempcount}}{#2}{#3}{#4}
}
\newcommand{\truefalseQuestion}[5]{\item<#1-|handout:#1-> \color<#2-|handout:#2->{#3} #4 \qquad \visible<#2-|handout:#2->{#5}}

\newcommand{\boder}{\ensuremath{\textcolor{blue}{\vee}}\text{} }
\newcommand{\bund}{\ensuremath{\textcolor{blue}{\wedge}}\text{}}
\newcommand{\bimp}{\ensuremath{\textcolor{blue}{\to}}\text{}}
\newcommand{\bnot}{\ensuremath{\textcolor{blue}{\neg}}\text{}}
\newcommand{\bone}{\ensuremath{\textcolor{blue}{1}}\text{}}
\newcommand{\bzero}{\ensuremath{\textcolor{blue}{0}}\text{}}
\newcommand{\bleftBr}{\ensuremath{\textcolor{blue}{(}}\text{}}
\newcommand{\brightBr}{\ensuremath{\textcolor{blue}{)}}\text{}}

\newcommand{\summe}[2]{\sum\limits_{#1}^{#2}}
\newcommand{\limes}[1]{\lim\limits_{#1}}

\newcommand{\numpp}{\advance \value{weeknum} by -2 \theweeknum \advance \value{weeknum} by 2}
\newcommand{\nump}{\advance \value{weeknum} by -1 \theweeknum \advance \value{weeknum} by 1}

\newcommand{\mycomment}[1]{}

\newcounter{abc}
\newenvironment{alist}{
  \begin{list}{(\alph{abc})}{
      \usecounter{abc}\setlength{\leftmargin}{8mm}\setlength{\labelsep}{2mm}
    }
}{\end{list}}

\renewcommand{\mod}{\text{ \bf{mod} }}
\renewcommand{\div}{\text{ \bf{div} }}

%\newcommand{\sgn}{\text{sgn}}


\graphicspath{{../figures/}}

\begin{document}
\starttut{6}

\lastframe{0.65}{25}{xkcd/advent_calendar.png}{https://www.xkcd.com/994/}


\begin{frame}{Zum letzten Übungsblatt}
	\only<beamer>{Korrektur dauert noch eine Woche.}
	\begin{itemize}
		\pause
		\item Vollständige Induktion: Die Induktionsvoraussetzung lautet\\
		{\enquote{für \textbf{ein} i gilt...}}.\\
		Für alle i zeigt ihr es erst in der Induktion!
		\pause
		\item Kaum jemand hat das richtig gemacht, es ist aber wirklich essentiell. Wer das bei der Klausur falsch macht, verschenkt einfache Punkte!
		\pause
		\item \color{myRed}Folgendes ist \textbf{nicht} korrekt, um A zu zeigen:\\
		$A \implies B$, B ist wahr. Also ist auch A wahr.\\
		Hier bedarf es Äquivalenz, damit aus B auch A folgt!
	\end{itemize}
\end{frame}

\framePrevEpisode

\begin{frame}{Rückblick: Codierungen}
	\begin{itemize}
		\item Codierungen: Injektive Abbildungen (meist Homomorphismen)
		\item $\varepsilon$-Freiheit, Präfixfreiheit
		\item Einfaches decodieren für präfixfreie Homomorphismen
		\item Huffman-Codierungen: Präfixfreie Codes mit optimaler Länge
	\end{itemize}
\end{frame}

\begin{frame}
	\begin{block}{Wahr oder falsch?}
	\begin{itemize}
		\truefalseQ{1}{2}{myGreen}{Das Zweierkomplement ist gut zum Rechnen}{W}
		\truefalseQ{3}{4}{myRed}{Jeder Homomorphismus ist präfixfrei}{F}
		\truefalseQ{5}{6}{myGreen}{Jede Huffman-Codierung ist präfixfrei}{W}
		\truefalseQ{7}{8}{myRed}{Jeder $\varepsilon$-freie Homomorphismus ist präfixfrei}{F $\qquad$ Aber:\\ Jeder präfixfreie Homomorphismus ist $\varepsilon$-frei}
		\truefalseQ{9}{10}{myGreen}{Präfixfreie Codes sind einfach zu decodieren.}{W}
	\end{itemize}
	\end{block}
\end{frame}

%% Zum Aufwärmen
\input{../Bloecke/Induktion_Vögel}

%% Inhalt

\section{Funktionen von Funktionen}

\begin{frame}{Funktionen}
	\begin{Definition}
		Seien A und B Mengen. Dann ist $$B^A := \set{f \Mid f \from A \functionto B }$$ die Menge aller Abbildungen von $A$ nach $B$.
	\end{Definition}

	Abbildungen kann man sich auch als Tabellen vorstellen: \\
	Wir wählen für jedes $a \in A$ ein $b \in B$.
	
	\begin{block}{Beobachtung}
		Für endliche Mengen A und B gilt:
		\delimitershortfall=0pt
		$$\setsize{B^A} = \setsize B^{\setsize A}$$
	\end{block}
\end{frame}

\newcommand{\filter}{filter_{\only<2-3|handout:1>{even}\only<4|handout:2>{odd}\only<5|handout:3>{Uglg}}}

\begin{frame}{Problem: Ziffern filtern}
	Wir wollen aus einer Ziffernfolge (Wort aus $Z^*_{10}$) Ziffern herausfiltern, die eine bestimmte Eigenschaft haben:\\
	\visible<2->{\impl\ }\only<2-3|handout:1>{Gerade}\only<4|handout:2>{Ungerade}\only<5|handout:3>{Löst die Ungleichung $x^2 - 3 < 20$}
	
	\visible<3-> {
		\begin{threealign}
		\filter \from Z^*_{10} &\functionto& Z^*_{10} \\
		\eps &\mapsto& \eps \\
		z \cdot v &\mapsto& \begin{cases}
		z \cdot \filter(v) &\text{falls } \only<2-3|handout:1>{z \in \{\word 2,\word 4,\word 6,\word 8,\word 0\}} \only<4|handout:2>{z \in \{\word 1,\word 3,\word 5,\word 7,\word 9\}} \only<5|handout:3>{\fnum_{10}(z)^2 - 3 < 20 }\\
		\filter(v) &\text{sonst}
		\end{cases}\\
		& \multicolumn{2}{l}{\text{wobei } z \in Z_{10} \text{ und } v \in Z_{10}^*.}
		\end{threealign}
	}
\end{frame}

\begin{frame}{Problem: Ziffern filtern}
	\begin{block}{Aufgabe}
		Schreibt die Funktionen zum Filtern aller Ziffern...
		\begin{itemize}
			\item $> 5$
			\item $< 5$
			\item $<2$ oder $>7$
			\item ... die perfekt sind
			\item ... die sich ausmalen lassen
		\end{itemize}
	\end{block}

	\pause
	\impl Blödsinn!\\
	Alle Funktionen unterscheiden sich nur \textbf{an einer Stelle}: \pause \\
	Der \textbf{Bedingung}, die überprüft wird.
\end{frame}

\begin{frame}{Filter}
	Definiere stattdessen eine \textbf{universelle} Filter-Funktion:
	
	\begin{threealign}
		filter_p : Z^*_{10} &\functionto& Z^*_{10} \\
		\eps &\mapsto& \eps \\
		z \cdot v &\mapsto& \begin{cases}
		z \cdot filter_p(v) &\text{falls } p(z) = \W\\
		filter_p(v) &\text{sonst}
		\end{cases}\\
		&\multicolumn{2}{l}{\text{wobei } z \in Z_{10} \text{ und } v \in Z_{10}^*}
	\end{threealign}
	\medskip
	und zwar für alle \quad $p: Z_{10} \functionto \BB$.
\end{frame}

\begin{frame}{Filter}
	\begin{threealign}
		filter_p : Z^*_{10} &\functionto& Z^*_{10} \\
		\eps &\mapsto& \eps \\
		z \cdot v &\mapsto& \begin{cases}
		z \cdot filter_p(v) &\text{falls } p(z) = \W\\
		filter_p(v) &\text{sonst}
		\end{cases}\\
		&\multicolumn{2}{l}{\text{wobei } z \in Z_{10} \text{ und } v \in Z_{10}^*}
	\end{threealign}
	
	Definiere beispielsweise nun:
	$$ even:  Z_{10} \functionto \BB, \; z \mapsto \begin{cases}
	\textbf{w} &\text{falls } z \in \{\word 2,\word  4,\word  6,\word  8,\word  0\}\\
	\textbf{f} &\text{sonst }
	\end{cases}$$
	
	\pause
	Dann ist $filter_{even}(\word{123456}) = \word{246}$.
	
	\smallskip
	\impl Brauchen bloß $p$-Funktionen definieren anstatt jedes Mal ganze Filter-Funktion!
	
\end{frame}

% TOO MUCH FOR THE FIRST TIME, I GUESS...
\mycomment{
	\begin{frame}{Charakteristische Funktion}
		Sei $M$ eine abzählbar unendliche Menge und $L \subseteq M$.
		
		\begin{Definition}
			Die \textbf{Charakteristische Funktion} einer (Teil-)Menge $L$ ist die Funktion 
			\begin{threealign}
			C_L : M &\functionto& \{0, 1\}\\
			x &\mapsto& \begin{cases}
			1 &x \in L \\
			0 &x \notin L
			\end{cases}.
			\end{threealign}
		\end{Definition}
		
		Also ist $C_L \in \{0, 1\}^M$.
		\smallskip
		
		Klar ist außerdem: Es gibt eine Bijektion $C$ zwischen Mengen und Charakteristischen Funktionen: $$C : 2^M \to \{0, 1\}^M, L \mapsto C_L$$
	\end{frame}
	
	\begin{frame}{Charakteristische Funktion}
		Jetzt können wir auf diesen Funktionen äquivalente Operationen wie auf Mengen definieren...
		
		Vereinigung: $V\colon \{0,1\}^M \times \{0,1\}^M\to \{0,1\}^M$\\[1em]
		
		\only<2|handout:1>{
			Beispielbild für $L_1=\{a,c,d\}$ und $L_2=\{b,c\}$
			
			\begin{tabular}{*{4}{>{$}c<{$}}}
				& L_1 & L_2 & L_1\cup L_2 \\
				x & f_1(x) & f_2(x) & V(f_1,f_2) \\
				a & 1 & 0 & 1 \\
				b & 0 & 1 & 1 \\
				c & 1 & 1 & 1 \\
				d & 1 & 0 & 1 \\
				e & 0 & 0 & 0 \\
			\end{tabular}
		}
		
		\only<3-|handout:2> {
			Wie definiert man $V(f_1,f_2)$? Zum Beispiel so:
			\begin{align*}
			V \colon \{0,1\}^M \times  \{0,1\}^M &\to \{0,1\}^M \\
			(f_1,f_2) &\mapsto (x \mapsto \max(f_1(x),f_2(x))) \\
			\end{align*}
		} \only<4|handout:2> {
		Oder so: $V(f_1,f_2) (x) = \max(f_1(x),f_2(x))$
	}
	\end{frame}
}

\newcommand{\yellow}[1]{\fcolorbox{white}{yellow!60}{\ensuremath{#1}}}
\begin{frame}{Funktionen auf LSD: Filter}
	\impl Idee: Mache $p$ zu einem Argument von $filter$!
	\medskip
	\begin{threealign}
		filter \from \yellow{\BB^{Z_{10}} \times } Z^*_{10} &\functionto& Z^*_{10} \\
		(\yellow{p}, \eps) &\mapsto& \eps \\
		(\yellow{p}, z \cdot v) &\mapsto& \begin{cases}
		z \cdot filter(\yellow{p, }v) &\text{falls } p(z) = \W\\
		filter(\yellow{p, }v) &\text{sonst}
		\end{cases}\\
		&\multicolumn{2}{l}{\text{wobei } z \in Z_{10} \text{ und } v \in Z_{10}^*.}
	\end{threealign} \\
	\pause
	
	\impl $filter(even, \word{123456}) = \word{246}$. \\
	\impl $filter(odd, \word{123456}) = \word{135}$.
	\medskip
	
	\impl $filter$ ist jetzt eine Funktion, die \textbf{eine Funktion als Argument} nimmt.
\end{frame}

\newcommand{\fconst}{\text{const}}
\begin{frame}{Funktionen auf LSD: const}
	Schnell: Gebt sieben verschiedene konstante Funktionen an. \\
	\pause
	\impl Zu aufwendig? \visible<6->{\impl Nicht mit \fconst: $\set{\fconst(1),\fconst(2),\fconst(3),...}$}
	\pause
	
	\begin{threealign}
		\fconst \from \R &\functionto& \R^\R \\
		\fconst(a) &:=& \underbrace{\big( x \mapsto a \big)}_{\text{eine Funktion!}}
	\end{threealign} \\
	\pause
	\impl Definiere $f := \fconst(42)$. \quad Dann ist $f(5) = f(1000) = f(-2.3) = 42$. \\
	\smallskip
	\pause
	Oder auch: $\left(\fconst(42)\right)(5) = 42$. \\
	\medskip
	\pause[7]
	
	\impl $\fconst$ ist eine Funktion, die \textbf{eine Funktion als Ergebnis} zurückliefert.
\end{frame}

\begin{frame}
	Wir haben gesehen:
	\begin{itemize}
		\item Eine Abbildung, die eine Funktion auf einen Wert abbildet
		\item Eine Abbildung, die einen Wert auf eine Funktion abbildet
	\end{itemize}
	Gleich werden wir das kombinieren! \\
	\medskip
	(Hinweis: Abbildung $=$ Funktion, gell? \smiley)
\end{frame}

\section{Speicher}

\begin{frame}{Bit vs. Byte}
	\begin{Definition}
		Eine \textbf{Bit} ist ein Zeichen des Alphabets $\{\word 0, \word 1\}$. \\
		Ein Wort aus 8 \emph{Bits} wird \textbf{Byte} genannt. 
	\end{Definition}
	\pause	
	
	\begin{Definition}
		Ein \textbf{Speicher} $$m \from \Adr \functionto \Val$$ bildet Adressen ($\Adr$) auf Werte ($\Val$) ab.\\
		Wir schreiben $\Mem$ für $\Val^\Adr$. \\
		Also: $m \in \Val^{\Adr}$ \quad bzw. \quad $m \in \Mem$.
	\end{Definition}
	\impl Betrachten abstrakte Funktionsweise anstatt konkrete Realisierung 
	%Hier interessiert uns nicht die konkrete Realisierung der Speicherung, sondern nur die abstrakte Funktionsweise.
	\pause
	
	\begin{block}{Hinweis}
		Häufig: Adressen und Werte als Binärzahlen. \\
		%Im Umgang mit Speichern benutzen wir hier nur Zahlen im Binärsystem. \\
		% Zu detailliert:                     und was ist mit negatiuven Zahlen?
		%$\Adr = \set{0, ..., 2^k-1}, \Val = {0,...,2^l-1} \quad \text{für } k,l \in \N_+$.
	\end{block}
\end{frame}

\begin{frame}{Speicher}

	Operationen: \\
	\begin{itemize}
		\item $\memread(m, adr)$, um eine Speicherzelle zu lesen 
		\pause
		\item $\memwrite(m, adr, val)$, um in eine Speicherzelle zu „schreiben“.  
	\end{itemize}
	
	\pause
	\begin{threealign}
		\memread \from \Val^{\Adr} \times \Adr &\functionto& \Val \\
		(m,a) &\mapsto& m(a)  \medskip  
		\visible<4->{ \\
			\memwrite \from \Val^{\Adr} \times \Adr \times \Val &\functionto& \Val^{\Adr} \\
			(m,a',v') &\mapsto& m' , \quad \text{wobei} 
		} 
		\visible<5->{ \\ 
				m'(a) &=&
				\begin{cases}
				v' & \text{ falls } a=a' \\
				m(a) & \text{ falls } a\not=a' \\
				\end{cases}	
		}
	\end{threealign}	
	\pause[6]
	$\memwrite$ gibt die neue Speichertabelle $m'$ zurück!
\end{frame}

\begin{frame}{Speicher}
	\begin{block}{Beispiel}
		
	$\memread(m, \word{01}) = \only<2->{\word{00000111}}$\\
	\visible<2-|handout:2>{$m' := \memwrite(m, \word{01}, \word{11111100})$\\}
	\visible<3-|handout:2>{$\memread\big(\underbrace{\memwrite(m, \word{01}, \word{11111100})}_{m'}, \word{01}\big) = \word{11111100}$\\}
	\bigskip
	
	\begin{table}
		\begin{tabular}{r}
			\\ \\ \\ \\  \\
			\visible<4-|handout:2>{Realworld-Analogie: }
		\end{tabular}
		\begin{tabular}{|c|c|}
			\hline
			\multicolumn{2}{|c|}{Speicher $m$} \\
			\hline
			\word{00} & \word{00101000} \\
			\word{01} & \word{00000111} \\
			\word{10} & \word{10010110} \\
			\word{11} & \word{00100101} \\ 
			\hline 
			\multicolumn{2}{c}{\visible<4-|handout:2>{Zustand bei $t=0$}}
		\end{tabular} 
		\visible<2-|handout:2>{
			\begin{tabular}{|c|c|}
				\hline
				\multicolumn{2}{|c|}{Speicher $m'$} \\
				\hline
				\word{00} & \word{00101000} \\
				\underline{\word{01}} & \underline{\word{11111100}} \\
				\word{10} & \word{10010110} \\
				\word{11} & \word{00100101} \\
				\hline
				\multicolumn{2}{c}{\visible<4-|handout:2>{Zustand bei $t=1$}}
			\end{tabular} 
		}
	\end{table}

	\end{block}
\end{frame}

\mycomment{
	\only<handout:0>{\begin{frame}{Aufgaben}
		
	\end{frame}}
}

\section{MIMA}

\begin{frame}{MIMA}
	\includegraphics[width=\linewidth]{MIMA_simple.png}\\
	Die \emph{MIMA} ist ein idealisierter Prozessor. 
\end{frame}

\begin{frame}{Eigenschaften}
	\begin{itemize}[<+->]
		\item Adressen sind 20 Bit lang
		\item \enquote{Werte} sind 24 Bit lang
		\item Befehlscodierungen -- zwei Formate:
		\begin{itemize}
			\item[a)] 4 Bit für den OpCode und 20 Bit für einen Parameter (Adresse / Konstante)
			\item[b)] 8 Bit Befehl (Rest irrelevant)\\
			\item[]<+(-2)-> \includegraphics[width=100px]{MIMA_commands.png} 
		\end{itemize} 
	\end{itemize}
\end{frame}


\begin{frame}{Wichtige Register}
	\begin{itemize}[<+->]
		\item \emph{IAR} : InstruktionsAdressRegister : Speichert Adresse des aktuell auszuführenden Befehls.
		\item \emph{IR}: InstruktionsRegister : Speichert den auszuführenden Befehl selbst.
		\item \emph{SAR}: SpeicherAdressRegister : Enthält die Adresse eines Wertes, der aus dem Speicher gelesen werden soll.
		\item \emph{SDR}: SpeicherDatenRegister : Enthält einen Wert, der aus dem Speicher geladen wurde.
		\item \emph{Akku}: enthält Ausgangswerte/Ergebnisse von Berechnungen
	\end{itemize}
\end{frame}

\begin{frame}[t]{Aufbau}
	\begin{figure}
		\centering
		\includegraphics[width=\linewidth]{MIMA.png}
	\end{figure}
\end{frame}

\begin{frame}[t]{Befehlsholphase}
	\begin{figure}
		\centering
		\includegraphics[width=\linewidth]{MIMA.png}
	\end{figure}
	\pause
	\only<2-6>{\begin{itemize}
		\only<2|handout:1> {\item[1.] IAR $\to $ SAR  und IAR $\to$ X 
			\item[] Befehlsadresse dem Speicher übergeben und Zähler zum Erhöhen an ALU geben.}
		\only<3|handout:2>{\item[2.] Eins $\to$ Y  
			\item[] 1-Wert für Erhöhung des Zählers an ALU geben.}
		\only<4|handout:3>{\item[3.] ALU aufaddieren ($Z=X+Y$) 
			\item[] Nächste Befehlsadresse berechnen. }
		\only<5|handout:4>{\item[4.] Z $\to$ IAR 
			\item[] Adresse für nächste Runde speichern. }
		\only<6|handout:5>{\item[5.] SDR $\to$ IR 
			\item[] Wert zur angefragten Adresse erhalten. }
	\end{itemize}}	
\end{frame}

\newcommand{\itemizeconfig}{\setlength{\parsep}{0pt}\setlength{\parskip}{0pt}\setlength{\topsep}{0pt}\setlength{\partopsep}{0pt}}
\newcommand{\explain}[1]{\hfill {\small #1 }}
\begin{frame}{Befehle}
	Die MIMA besitzt einen Befehlssatz mit möglichen Befehlen. %Andere Befehle (oder Varianten) werden NICHT unterstützt und können daher nicht verwendet werden!
	\textbf{Nur die Befehle aus der VL} dürft ihr benutzen! \\ 
	\smallskip
	(Befehle, die eigentlich zwei Operanden brauchen, lesen zusätzl. vom Akku.)
	{
	\begin{itemize}[<+->] \itemizeconfig
		\item Rechenoperationen:
		\begin{itemize} 
			\item ADD adr \explain{Akku $\leftarrow$ Akku + Mem(adr)}
			\item AND, OR, XOR adr \explain{Dasselbe mit bitweisem Und, Oder und eXclusive OR}
			\item NOT, RAR \explain{Akku $\leftarrow$ not(Akku); \quad  Rotate-Akku-Right();}
		\end{itemize}
		\item Datentransport:
		\begin{itemize}
			\item LDC const \explain{Akku $\leftarrow$ const\quad (const ist dabei eine \textbf{20-Bit}-Konstante) }
			\item LDV, STV adr \explain{(Lese/Schreibe in/vom Akku von/an Adresse)}
			\item LDIV, STIV adr \explain{(Lese/Schreibe von/an Adresse, die an der Adresse steht)}
		\end{itemize}
		\item Vergleichsoperation: EQL adr \explain{(liefert $-1$ wenn Akku $=$ Mem(adr), 0 sonst)}
		\item Sprünge:
		\begin{itemize}
			\item JMP adr \explain{Springe zu Befehl in Adresse adr}
			\item JMN adr \explain{Springe, falls Akku negativ \quad „JuMp if Negative“}
		\end{itemize}
		\item HALT \explain{Stoppt die MIMA}
	\end{itemize}}
\end{frame}

\begin{frame}{Bemerkungen}
	\begin{block}{Indirekte Adressierung (LDIV, STIV)}
		\centering
		\includegraphics[width=150px]{MIMA_indirect.png}
	\end{block}
	
	\pause
	\begin{block}{HALT}
		Jedes Programm muss mit HALT enden! Sonst läuft das Programm endlos weiter! (Kostet dann Punkte!)
	\end{block}

	\pause
	\begin{block}{Negative Konstanten}
		Negative Konstanten können \textbf{nicht} mit LDC geladen werden. Warum? \\
		\pause \impl Unser Akku ist 24 Bit breit, aber wir können nur in die \textbf{hinteren 20 Bit} laden!
	\end{block}
\end{frame}

\begin{frame}{Negative Zahlen}
	\begin{block}{Aufgabe}
		Schreibt ein Programm, das von einer an Adresse $a_1$ gegebenen positiven Zahl das Zweierkomplement berechnet und an Adresse $a_2$ ablegt.
	\end{block}

	\visible<2|handout:2>{
		\begin{block}{Lösung}
			LDV $a_1$\\
			NOT\\
			STV $a_2$\\
			LDC 1\\
			ADD $a_2$\\
			STV $a_2$
		\end{block}
	}
\end{frame}

\begin{frame}{Beispiele}
	Beispiele zur Umsetzung von Anweisungen aus Hochsprachen (if, while, for....):\\
	Siehe Übung WS 15/16
\end{frame}

\begin{frame}{Übung: Modulo 2 und Betrag}
	1. Schreibe ein Programm, das eine an Speicheradresse $a_1$ gegebenen Zahl Modulo 2 rechnet und an Adresse $a_2$ ablegt. \\
	\medskip
	2. Schreibe ein Programm, das den Betrag einer an Adresse $a_1$ gegebenen Zahl berechnet und an Adresse $a_2$ ablegt.
\end{frame}

\begin{frame}{Lösung: Modulo 2}
	\begin{tabbing}
		start: \; \= LDC 1 \quad  //$= 000000000000000000000001_2$ \\
				\> AND $a_1$ \\
				\> STV $a_2$ \\
				\> HALT
	\end{tabbing}
\end{frame}

\begin{frame}{Lösung: Betrag}
	\begin{tabbing}
		start: \quad \= LDV $a_1$ \\
					 \> JMN negate \\
		end: 		 \> STV $a_2$ \\
					 \> HALT \\
		negate:		 \> NOT \\
					 \> STV $a_2$ \\
					 \> LDC 1 \\
					 \> ADD $a_2$  \\
					 \> JMP end	 \\		 
	\end{tabbing}
\end{frame}

\begin{frame}{Übung: Modulo 3}
	Schreibe ein Programm, das eine an Speicheradresse $a_1$ gegebenen Zahl Modulo \textbf{3} (drei!) rechnet und an Adresse $a_2$ ablegt.
\end{frame}

\begin{frame}{Lösung: Modulo 3}
  \begin{tabbing}
    start: \quad \= LDC 1 \\
           \> STV One\\
           \> LDC 3\\
           \> NOT\\
           \> ADD One\\
           \> STV MinusThree\\
           \> LDV $a_1$\\
    \medskip
    while: \> ADD MinusThree\\
           \> JMN end\\
           \> JMP while\\
    \medskip
    end:   \> STV $a_1$\\
           \> LDC 3\\
           \> ADD $a_1$\\
           \> STV $a_2$ \\
           \> HALT\\
  \end{tabbing}
\end{frame}


\begin{frame}	
	\begin{block}{Was ihr nun wissen solltet}
		\begin{itemize}
			\item Speicher - Abbildungen, die Abbildungen auf Abbildungen abbilden!
			\item Wie ein einfacher Prozessor aufgebaut ist
			\item Wie man die MIMA für einfache Aufgaben nutzt
		\end{itemize}
	\end{block}
	
	\begin{block}{Was nächstes Mal kommt}
		\begin{itemize}
			\item Grammatik ist schwer - aber nicht bei uns!
		\end{itemize}
	\end{block}
\end{frame}


\lastframe{0.55}{30}{xkcd/processor.png}{https://www.xkcd.com/676/}
\slideThanks

\end{document}